% !TEX root = ../main.tex
% !TEX program = lualatex
\documentclass[../main.tex]{subfiles}
\IfSubfilesClassLoaded{\externaldocument{\subfix{../build/main}}}{}
% =============================================================================
% BOARD ASSUMPTIONS AND FAST CUES
% DESCRIPTION: Overarching assumptions plus fast cues/common board pivots.
% =============================================================================

\begin{document}
\IfSubfilesClassLoaded{\chapter{Board Assumptions and Fast Cues}}{}
%====================
% Contents here
%====================

% === Assumptions ===
\section{Assumptions}

\subsection{Summary}
\begin{itemize}
  \item This chapter is the board-speed crosswalk for unstable facts, recurring traps, and high-yield pivots. Treat it as a final rehearsal checklist, not as a replacement for the detailed chapter sources.
  \item The safest oral-board pattern is: identify the authority, name the current source or caveat, separate adjacent concepts, and then walk the consequence for cost, schedule, performance, technical risk, or Fleet readiness.
  \item The most source-sensitive facts are flag billets, PAE/PEO roles, statutory thresholds, FAR/DFARS/NMCARS text, committee names, and current program categorizations. Verify them close to the board date.
  \item Weak-area cues should rotate across relationships and exact-list recall: SE baselines, C4I product ownership, APB vs.\ PMB, requirements-to-acquisition flow, fiscal law, ship-change execution, CIVPERS routes, and current EDO organization.
\end{itemize}

\begin{description}
  \item[\textbf{Currency.}] Flag assignments, \ac{peo} names, \ac{frd} numbering, stabilized rates, and statutory thresholds change mid-year; confirm figures and approval levels in the week of use. \ac{far}/\ac{dfars}/\ac{nmcars} revisions release frequently, so treat citations as baseline until superseded~\autocite{FAR,DFARS,NMCARS}.
  \item[\textbf{Governance anchors.}] \ac{don} organizational relationships, \ac{ta}/\ac{pa} roles, and delegated authorities follow~\autocite{SECNAV5400-15D,SECNAV5000-2G}; pathway policy and cost-estimate governance per~\autocite{DoDI5000-85}.
  \item[\textbf{Execution anchors.}] Definitions of commitment, obligation, expenditure/expense, and outlay follow~\autocite{DoDFMR-Vol3}; \ac{igt} flow uses Treasury G-Invoicing~\autocite{Treasury-GInvoicing}.
  \item[\textbf{Fiscal law.}] Purpose (\S~1301), Bona Fide Need (\S~1502), and Anti-Deficiency (\S~1341/\S~1517) statutes frame every vignette~\autocite{USC-1301-Purpose,USC-1502-BFN,USC-1341-ADA}.
  \item[\textbf{Rates \& NWCF.}] \ac{slr}, \ac{nor}, \ac{aor}, and \ac{nwcf} vs.\ mission-funded distinctions derive from the latest \ac{fmr} updates~\autocite{DoDFMR-Vol3,DoDFMR-Vol3-Ch19}.
  \item[\textbf{Contracting overlays.}] \ac{ko} authority, cost principles, source-selection doctrine, and shipbuilding execution rely on the current \ac{far}/\ac{dfars}/\ac{nmcars} stack plus active \ac{rfo} class deviations and \ac{navsea} implementation guides. Monitor \ac{navsea} notices for \ac{som}, Source Selection Guide, and Contracts Handbook updates; public access checks of the \ac{navsea} guide endpoints returned 403 on 2026-06-05, so use them as retained official baselines until manually verified~\autocite{FAR,DFARS,NMCARS,FAR-Overhaul-Deviation-Guide-2026,SOM-Ch3-2023,NAVSEA-Source-Selection-Guide-2022,NAVSEA-Contracts-Handbook-2023}.
  \item[\textbf{Competition rules.}] \ac{cica} remains the statute anchor for full-and-open competition, but current DoD practice uses the active \ac{rfo} Part~6/DFARS~206 deviation. Treat the seven authorities as a memory hook, then verify revised \ac{far}~6.103/\ac{far}~6.104/\ac{far}~6.301 before quoting legacy 6.302/6.303/6.305 labels~\autocite{USC-3201-CICA,FAR-Overhaul-Part6-2026,DoD-RFO-Deviation-Part6-2026}.
  \item[\textbf{Solicitation prep.}] Release criteria, evaluation governance, and communication rules follow \ac{far} parts~5/15, the \ac{navsea} Source Selection Guide/Contracts Handbook, and Coursebook~3.2.2; treat as baseline until updated~\autocite{FAR,NAVSEA-Source-Selection-Guide-2022,NAVSEA-Contracts-Handbook-2023,edo-3-2-2-solicitation-prep-2025}.
  \item[\textbf{Coursebook crosswalk.}] Content tracks \ac{edo} courseware modules cited in the bibliography (2.1.*, 3.1.*, 3.2.*) for board prep currency~\autocite{edo-2-1-1-command-structures-2025,edo-2-1-5-technical-authority-2025,edo-3-1-1-ppbe-2025,edo-3-1-4-budget-execution-2025,edo-3-2-1-intro-contracting-2025,edo-3-2-2-solicitation-prep-2025}.
\end{description}
\HR

% === STUDY CUES ===
\section{Fast Cues and Common Board Pivots}
\subsection{High-Yield Rotating Weak-Area Cues}
\begin{description}
  \item[\textbf{Requirements-to-acquisition chain.}] Can you walk from capability gap to \ac{icd}, \ac{mdd}, \ac{msa}/\ac{aoa}, \ac{tmrr}, validated \ac{cdd}, Milestone~B, and \ac{emd} without mixing who validates requirements with who authorizes acquisition phase entry~\autocite{CJCSM5123-01,DoDI5000-85,edo-3-3-2-acq-milestones-phases-2025,edo-3-3-3-jcids-2025}?
  \item[\textbf{APB vs.\ PMB.}] \ac{apb} is the MDA-approved program cost/schedule/performance baseline; \ac{pmb} is the contractor \ac{evm} execution baseline. Bad CPI/SPI starts with variance/root-cause analysis before an \ac{apb} revision discussion~\autocite{DoDI5000-85,edo-3-6-1-evm-intro-2025,edo-3-6-2-evm-analysis-2025}.
  \item[\textbf{Technical-review trap.}] \ac{sfr} supports the functional baseline, \ac{pdr} supports the allocated baseline, and \ac{cdr} supports the initial product baseline; \ac{srr} reviews requirements but normally is not the baseline-setting review~\autocite{DoDI5000-88,DAU-Technical-Baselines,DAU-System-Requirements-Review,DAU-System-Functional-Review,DAU-Preliminary-Design-Review,DAU-Critical-Design-Review}.
  \item[\textbf{KPP/KSA/APA and test evidence.}] If a board question says ``not a \ac{kpp},'' still ask whether the missed attribute is a \ac{ksa} or \ac{apa}, whether it is measurable/testable, and what the \ac{apb}, \ac{temp}, \ac{dt}, \ac{ot}, and residual-risk consequences are~\autocite{CJCSM5123-01,DoDI5000-85,edo-3-3-3-jcids-2025}.
  \item[\textbf{EDO flag count.}] Current \ac{edo} flag roster recall is 3~\acp{vadm}, 4~\acp{radm}, 1~RADM(sel), and 7~\acp{rdml}. Count Andrew Biehn as a PAE Maritime / NAVSEA 21 awareness cue, not an \ac{edo} flag-roster entry~\autocite{navy-flag-roster-2026-06,PAE-Maritime-Homepage-2026}.
  \item[\textbf{T\&E and ship-trials split.}] \ac{dte}/\ac{techeval} builds technical confidence; \ac{optest}/\ac{opeval}/\ac{iote} proves operational effectiveness and suitability through the \ac{ota}; \ac{lfte} covers survivability/lethality for covered systems; and \ac{insurv}/ship trials support material readiness, acceptance, and delivery decisions. Alternative-pathway \ac{te} is designation-specific, not the default answer~\autocite{USC-10-4171,USC-10-4172,USC-10-8674,DoDI5000-98,OPNAVINST-4700-8M}.
  \item[\textbf{C4I product ownership.}] Separate combat-system C2/sensor netting (\ac{iws} 6.0/CEC), operational C2 applications (PMW~150/GCCS-M/MTC2/C2P), shipboard network and routing infrastructure (PMW~160/CANES/ADNS), communications/nav (PMW/A~170), tactical data-link terminals (PMA/PMW~101), and expeditionary C2 (PMW~790/DJC2)~\autocite{navsea-cec-enhancing-battlefield-awareness-2025,PEOC4I-PMW150-tearsheet-2026,PEOC4I-PMW160-tearsheet-2025,PEOC4I-PMW170-tearsheet-2025,PEOC4I-PMA-PMW101-tearsheet-2025,PEOC4I-PMW790-tearsheet-2025}.
  \item[\textbf{CIVPERS route choice.}] For poor performance, first ask whether the employee is probationary/trial or tenured. Tenured Chapter~43 means \ac{pip}/opportunity period plus substantial evidence; Chapter~75 means no required \ac{pip}, but preponderance, efficiency-of-service nexus, due process, and Douglas Factor penalty analysis. \ac{wg}/\ac{wl}/\ac{ws} are Federal Wage System nonsupervisory/leader/supervisory cues: 15/15/19 grades~\autocite{opm-performance-management-federal-employees-2026,opm-probationary-periods-2025,opm-employee-rights-appeals,ecfr-5cfr-532-203}.
  \item[\textbf{Modernization-to-fleet chain.}] Can you walk from \ac{pmo}/\ac{isea} need through \ac{scd}/\ac{ndenm} approval, installation package maturity, \ac{rmmco} gatekeeping, \ac{ait}/\ac{nsa}/\ac{lma} execution, \ac{sovt}, \ac{ils} handover, and \ac{cdmdoa}/as-installed closeout without skipping who owns the next risk decision~\autocite{JFMM-VolVI-RevD-Chg5-2025,TS9090-310,NAVSEA-Logistics-Planning-Processes-Policy-2026,edo-4-1-2-navy-modernization-2025,edo-4-1-4-ait-fundamentals-2025}?
  \item[\textbf{Maintenance closeout.}] For ship changes and availabilities, distinguish work package maturity, production completion, signed test evidence, \ac{sovt} closeout, lessons learned, and who owns the next risk decision~\autocite{COMFLTFORCOMINST4790-3,TS9090-310,edo-4-2-1-cno-availability-planning-2025,edo-4-2-2-cno-availability-execution-2025}.
  \item[\textbf{Manufacturing readiness.}] Before Milestone~C or rate decisions, ask whether design, suppliers, tooling, facilities, workforce, quality controls, process capability, and sustainment are mature enough to produce at rate; quality is built into the process, not inspected into the product~\autocite{edo-3-7-1-manufacturing-producibility-2025,edo-3-7-2-production-quality-manufacturing-2026,DoD-MRL-Deskbook-2025,DoD-Producibility-Manufacturability-Guide-2024}.
\end{description}

\subsection{Contracting, Fiscal, and Cost Pivots}
\begin{description}
  \item[\textbf{Authority split.}] Who owns \textbf{\ac{pa}} vs.\ \textbf{\ac{ta}}? Who can \textbf{bind the Government}, and within what delegated \ac{ko} authority~\autocite{SECNAV5400-15D,SECNAV5000-2G,FAR}?
  \item[\textbf{PM/KO huddle.}] Before release or modification, did the \ac{pm} and \ac{ko} align on requirements stability, acquisition strategy, and competition decisions~\autocite{DoDI5000-85,edo-3-2-1-intro-contracting-2025}?
  \item[\textbf{Release checklist.}] Synopsis filed, funds certified, SSP/\ac{oci} paperwork signed, and Sections~L/M synchronized before dropping the \ac{rfp}~\autocite{FAR,NAVSEA-Source-Selection-Guide-2022,edo-3-2-2-solicitation-prep-2025}?
  \item[\textbf{Spec departure.}] Which \textbf{\ac{twh}} warrants apply? What's the path for \emph{minor} vs.\ \emph{major} (safety-critical) deviations~\autocite{SECNAV5400-15D}?
  \item[\textbf{Competition trigger.}] Does \ac{cica} demand full and open? If not, which revised FAR~6.103 authority applies, and what level approves the \acs{ja} under revised FAR~6.104/DFARS~206.104~\autocite{USC-3201-CICA,FAR-Overhaul-Part6-2026,DoD-RFO-Deviation-Part6-2026}?
  \item[\textbf{JA vs.\ DF.}] \ac{ja} or \ac{df}? Know the revised approval levels, special cases such as public interest, and posting rules now organized under revised FAR~6.104 and 6.301~\autocite{FAR-Overhaul-Part6-2026,DoD-RFO-Deviation-Part6-2026}.
  \item[\textbf{Responsiveness vs.\ responsibility.}] For sealed bids, did the offer meet every requirement? Does the awardee meet the responsibility standards in~\autocite[\S~9.104]{FAR}?
  \item[\textbf{Method choice.}] Can the requirement be sealed bid (price-only, no discussions) or do risk/innovation drivers force negotiated procurement~\autocite{FAR,NAVSEA-Source-Selection-Guide-2022}?
  \item[\textbf{Contract type risk.}] Who owns cost risk (\ac{ffp}/\ac{fpif} vs.\ \ac{cpff}/\ac{cpaf}/\ac{cpif}) and which \ac{mca} phase justifies the selection~\autocite{FAR,DoDI5000-85}?
  \item[\textbf{FPIF math.}] Ready to compute share ratio, contract price, and \ac{pta} from target cost/profit and ceiling price if the Board throws numbers at you~\autocite{FAR,NAVSEA-Contracts-Handbook-2023}?
  \item[\textbf{Flexible forms.}] Do you know when to reach for \ac{idiq}, \ac{tnm}/\ac{lh}, or a \ac{uca}---and the surveillance and definitization rules each triggers~\autocite{FAR,NMCARS,DFARS}?
  \item[\textbf{Simplified acquisition.}] Above or below the current \$350K \ac{sat}? Are you maximizing \ac{sap} benefits while meeting small-business reservation rules~\autocite{FAR,AcquisitionGov-Threshold-Changes-2025}?
  \item[\textbf{Small business levers.}] Which \ac{sba} program (8(a), \ac{hubzone}, \ac{sdvosb}, \ac{wosb}) best fits the requirement, \ac{naics} size standard, rule-of-two evidence, and documentation trail~\autocite{FAR-19-102-Size-Standards,FAR-19-502-2-Set-Asides,SBA-Size-Standards-2026}?
  \item[\textbf{Innovation tools.}] Is this better suited for a 10~U.S.C.\ \S~4022 \ac{ot} prototype/follow-on, a 10~U.S.C.\ \S~4021 research \ac{ot}, or an \ac{sbir}/\ac{sttr} Phase I/II/III award path with data rights and transition protection~\autocite{USC-4021-OTA,USC-4022-OTF,USC-638-SBIR,SBIR-Apply-2026}?
  \item[\textbf{Data rights.}] Do you understand whether deliverables are unlimited, government-purpose, or limited/restricted and how that shapes future competition~\autocite{DFARS}?
  \item[\textbf{UCF alignment.}] Are Sections~A through H, I through K, and L and M synchronized, with factors and instructions perfectly mirrored~\autocite{FAR,NAVSEA-Source-Selection-Guide-2022}?
  \item[\textbf{Debrief/protest clock.}] For \ac{dod} source selections, do not stop at the \ac{far} three-day request/five-day debrief memory hook; check \ac{dfars} enhanced-debriefing follow-up questions and the \ac{gao} stay clock before advising a protest response~\autocite{FAR,DFARS}.
  \item[\textbf{Order type.}] Is this \textbf{\ac{igt}} (7600A/B) or a \textbf{contract}? When do we obligate, accept, and pay~\autocite{DoDFMR-Vol3,Treasury-GInvoicing}?
  \item[\textbf{Color of money \& window.}] Why \ac{rdte} vs.\ \ac{opn} vs.\ \ac{omn}? What's the obligation/expense horizon and expense rate~\autocite{DoDFMR-Vol3}?
  \item[\textbf{NWCF vs.\ Mission-Funded.}] \ac{slr} in play? How does that shift \textbf{\ac{nor}}/\textbf{\ac{aor}}? Which billets stay mission-funded~\autocite{DoDFMR-Vol3}?
  \item[\textbf{Reprogramming ``fork''.}] Map \textbf{\ac{btr}}, \textbf{\ac{ir}}, \textbf{\ac{pa-reprog}}, and \textbf{\ac{lt}} to the threshold or color-of-money change that forces each action~\autocite{DoDFMR-Vol3}.
  \item[\textbf{Constructive change vs.\ unauthorized commitment.}] Did Government conduct change performance, or did someone without authority make an agreement? Constructive changes drive \ac{ko} modification/equitable-adjustment or claim analysis; unauthorized commitments require FAR~1.602-3 ratification review before the Government is bound~\autocite{FAR,SOM-Ch3-2023,NAVSEA-Contracts-Handbook-2023}.
  \item[\textbf{Fiscal law flashcards.}] Decide whether \textbf{Purpose} (31~U.S.C.~\S~1301), \textbf{Bona Fide Need} (31~U.S.C.~\S~1502), or \textbf{\ac{ada}} (31~U.S.C.~\S~1341/\S~1517) governs the vignette presented~\autocite{USC-1301-Purpose,USC-1502-BFN,USC-1341-ADA}.
  \item[\textbf{Cost talk.}] Direct vs.\ indirect, pick an estimating method (analogy/parametric/engineering/actuals), call out learning-curve or escalation assumptions~\autocite{DAU-AAF,FAR}.
  \item[\textbf{Lifecycle framing.}] Distinguish Flyaway, Weapon-System, Procurement, and \textbf{\ac{lcc}}/\textbf{\ac{toc}} values so you can explain exactly what you are quoting and why~\autocite{DoDI5000-85,SECNAV5000-2G}.
  \item[\textbf{OCI / independence cue.}] Why is an organic Warfare Center (can't/shouldn't/won't) preferred vs.\ industry for this scenario~\autocite{SECNAV5400-15D}?
\end{description}

\subsection{Quick Review}
\begin{enumerate}
  \item Pick any cue and state the governing authority first: statute, policy, instruction, source-selection rule, or current official organization page.
  \item For APB vs.\ PMB, SE review baselines, and DT/OT/LFT\&E, explain the difference without using similar-sounding words as substitutes.
  \item For a C4I product question, map the product to the correct office before discussing ACAT or pathway status.
  \item For a fiscal-law vignette, identify purpose, time, and amount in that order before choosing a color of money or corrective path.
  \item For a ship-change or availability vignette, walk from program-office need to ISEA/TA review, installation package, RMMCO/NSA/LMA execution, SOVT, ILS/configuration handover, and risk owner.
\end{enumerate}
%====================
% End of file
%====================
\IfSubfilesClassLoaded{
  \RenewDocumentCommand{\entryname}{}{\textbf{\color{Modern} Acronym}}
  \RenewDocumentCommand{\descriptionname}{}{\textbf{\color{Modern} Definition}}
  \printnoidxglossary[
    type=\acronymtype,
    title=Acronyms,
    style=long-booktabs
  ]}{}
\end{document}
